\subsubsection{Source databases and reference compatibility}
\label{sec:datasets}
Table~\ref{tab:datasets} distinguishes the source families retained in historical development from the sources used for later analyses. A source name identifies an acquisition or repository family; a source alias in a derived file does not necessarily identify a new hospital or an independent patient population. PTB-XL appears both as native SCP-coded data and through Challenge-format aliases. Similarly, a derived sampling-rate version does not create a new cohort. The source descriptions below therefore concern reference provenance and intended analytical use, not a claim that all historical source partitions have been authenticated.

PTB-XL comprises ten-second 12-lead ECGs with hierarchical diagnostic, rhythm and form statements and source patient identifiers \citep{Wagner2020PTBXL}. Its published patient-based folds provide a useful contemporary partition convention, but observing a fold identifier in a recovered evaluation table does not establish which records every preceding training job consumed. The P4 evaluation uses selected positive and negative panels derived from SCP statements. It should not be interpreted as a representative sample of infarction, hypertrophy or ischemia prevalence. The four panels overlap at the ECG level, and a known target in one panel does not imply that all other pathology targets are known for that ECG.

Chapman--Shaoxing is a Chinese 12-lead arrhythmia database described by Zheng et al. \citep{Zheng2020Chapman}. It dominates the retained P6 sample. CPSC2018, CPSC2018-Extra, Georgia and Ningbo are retained as distinct source labels in the Challenge-associated files \citep{Reyna2022Challenge}. Their diagnoses are mapped to a shared output vocabulary while their original annotation procedures remain source-specific. CODE test is the public expert-reference test resource associated with the Brazilian 12-lead study of Ribeiro et al. \citep{Ribeiro2020Automatic}; the 19 selected P6 records from that resource have already contributed to the reported development evidence. They cannot subsequently serve as a newly untouched external cohort for this same performance claim.

SPH denotes \emph{Shandong Provincial Hospital}. Its large multi-label database provides standardized diagnostic statements \citep{SPH2022}. The completed retrospective experiment used 25,770 ECGs and 24,666 source patient identifiers. Its mapping supplies five P6, six P8 and three P4 compatible endpoints. A source Normal-only statement is not equivalent to a positive NSR rhythm annotation. Absence of a mapped diagnostic statement is an absent-statement comparison within this mapping; it is not an independently adjudicated disease-negative control. These distinctions are particularly important for rare pathology endpoints because a large number of absent statements can strongly alter predictive values.

LUDB provides 200 ten-second 12-lead recordings with waveform annotations and diagnostic metadata \citep{LUDB2021,LUDB2020Paper}. The retrospective route study used all 200 ECGs. Explicit principal rhythm contrasts support some endpoints, whereas incidental findings may provide positives without reliable negative references. For example, LVH-positive and unknown LUDB records must not be combined into a conventional positive/negative P4 panel. PTB Diagnostic is a separate source containing 549 ECGs from 290 subjects \citep{PTBDiagnostic}. The pathology transport analysis used named infarct territories as positive references and explicitly healthy controls; other records were retained as unknown for each territory comparison. Neither LUDB nor PTB Diagnostic was treated as a replacement for the missing parent population of the June specialist panels.

\begin{table}[tbp]
\caption{Source families, reference types and analytical roles. Counts are study-specific observations, not a reconstructed joint source-by-partition roster.\label{tab:datasets}}
\centering\small
\begin{tabular}{@{}p{.19\linewidth}p{.25\linewidth}p{.26\linewidth}p{.23\linewidth}@{}}
\toprule
Source / origin & Reference basis & Use in this study & Boundary \\
\midrule
PTB-XL / Germany & SCP diagnostic, rhythm and form statements & P4 selected panels; historical development; later folds 1--9 bridge & Native and Challenge aliases can overlap \\
Chapman--Shaoxing / China & Source rhythm diagnoses & 1,514 selected P6 ECGs & Dominant retained source; historical patient isolation unverified \\
CPSC2018, Extra; Ningbo / China & Challenge-associated diagnostic codes & Historical development and selected P6 aliases & Source-specific annotation; small selected strata \\
Georgia / USA & Challenge-associated diagnostic codes & Historical development and six selected P6 ECGs & Small retained evaluation stratum \\
CODE test / Brazil & Supplied expert-reference labels & 19 selected P6 ECGs & Already inspected; no new independent-test claim \\
Mixed local group & Retained heterogeneous reference fields & 63 selected P6 ECGs & Not one institution or a uniform adjudication protocol \\
SPH / China & Standardized multi-label statements & Full 25,770-ECG retrospective analysis & Missing 1AVB, NSR and ISC\_ mappings \\
LUDB / Russia & Principal rhythm and diagnostic metadata & 200-ECG route and reconstruction studies & Unknown references retained as unknown \\
PTB Diagnostic / Germany & Subject diagnosis, named territory, healthy controls & 549-ECG pathology transport analysis & Endpoint-specific known sets; unknowns excluded only from endpoint metrics \\
\bottomrule
\end{tabular}
\end{table}

\subsubsection{Partition accounting and exclusions}
\label{sec:partition-accounting}
Table~\ref{tab:historical-partitions} records recoverable population counts without combining incompatible units. A file can contain several windows; a patient can contribute several files; and one ECG can contribute to more than one binary pathology panel. Consequently, the total number of training patients cannot be obtained by adding task-specific window counts. The source counts of the archived job 16768 sum to its 9,945 windows: CPSC2018, 3,394; CPSC2018-Extra, 45; Georgia, 659; INCART, 359; Ningbo, 22; PTB-XL, 5,420; PTB-XL Challenge 2020, three; extra 3AVB, 29; and Ningbo 3AVB, 14. These are source-attributed windows for that run, not a release-wide source-by-class training census.

\begin{table}[tbp]
\caption{Recoverable historical development and selected-evaluation counts. NR means not recovered with sufficient job-level binding.\label{tab:historical-partitions}}
\centering\small
\begin{tabular}{@{}p{.26\linewidth}p{.23\linewidth}p{.23\linewidth}p{.21\linewidth}@{}}
\toprule
Stage & Training & Validation / calibration & Statistical unit and limit \\
\midrule
V3C / job 16768 & 9,021 & 924 validation & Windows; 7,485 enumerated files; 390 record matches removed in roster construction \\
V3F / job 16781, rhythm & 9,021 & 144 validation & Task-specific windows; eight reported epochs \\
V3F / job 16781, pathology & 5,013 & 776 validation & Task-specific windows, not distinct patients \\
V3M / V3AG split & 9,885 & 1,132 validation & Record-based split; possible patient overlap in conditional reconstruction \\
P4 calibration / job 17296 & NR & 1,492 & Window--panel observations; exact consumed roster NR \\
Released P6 selected profile & Not inferred & 1,643 evaluation ECGs & Development and model selection \\
Released P8 profile & Not inferred & 2,193 evaluation ECGs & Development validation; overlaps P6 \\
Released P4 panels & Not inferred & 962 panel observations & 733 distinct ECGs; not 962 patients \\
\bottomrule
\end{tabular}
\end{table}

Known exclusions must be distinguished from unverifiable exclusions. The V3C record-match removal is documented, but it does not certify removal of every patient or signal duplicate from all subsequent stages. The recovered V3M/V3AG record split permits a conditional reconstruction with possible shared patients. It does not identify the exact bytes consumed by the historical jobs, and it must not be converted into a claim that the official PTB-XL fold-10 test partition was contaminated. Conversely, absence of a demonstrated collision in a bounded signature search is not proof that all historical data were independent. Missing consumed rosters and incomplete checkpoint ancestry prevent a definitive historical independence statement.

Window preparation follows the fixed tensor contract; technically invalid or incompatible references are tracked separately from disease-negative labels. The selected evaluation files do not retain a complete preselection population with every exclusion reason. Therefore, a June all-input enrollment diagram, source-by-class-by-patient partition table or natural-prevalence clinical predictive value cannot be reconstructed from the selected records alone. Later all-input studies preserve their own enrolled denominators and technical failures. Their eligibility accounting is reported with those studies rather than retroactively attached to the original June analysis.

For the historical partitions, we use three evidence states. \emph{Recorded} denotes surviving job summaries, loss/launcher excerpts or archived selected prediction rows: they support the quantities they contain, including windows or ECG--label observations, without proving the identities actually consumed by every trainer. \emph{Reconstructed} denotes later joins of source metadata, aliases, folds and record identifiers; these permit conditional overlap checks but cannot certify the missing job inputs. \emph{Uncertified} denotes the unrecovered joint source-by-class-by-partition patient roster, job-bound label masks, input checkpoint ancestry and complete historical patient/signal separation. Uncertified cells remain unavailable, rather than zero or estimated. The contemporary within-CODE audit has its own sealed split evidence (Section~\ref{sec:code-methods}) and does not repair this historical gap. Supplementary file \path{PARTITION_EVIDENCE_STATUS.csv} records these distinctions and their source files.

\subsubsection{Independent external evaluation in progress}
\label{sec:external-plan}
The purpose of the ongoing independent study is to test frozen ECG-statement prediction beyond the inspected development analyses. Independence is a property of the complete development and evaluation process, not of a database name. A candidate cohort can be considered confirmatory only after its admissible labels, patient and signal overlap checks, technical eligibility and prior inspection history have been documented. Already inspected SPH data, development-fold PTB-XL data, selected CODE test records and data used for fitting residual heads or adapting thresholds do not qualify automatically.

The completed retrospective analyses establish executable preprocessing and scoring contracts for specified checkpoints, but do not by themselves establish an admissible independent cohort. Before any independent performance estimate is incorporated, the analysis must bind the cohort roster and source version, label mapping and unknown-reference rules, checkpoint filenames and SHA-256 digests, inference configuration, score aggregation, membership cutoffs and statistical plan. Both source identifiers and waveform-duplicate checks are required because aliases or repeat acquisitions can evade record-name matching. Unresolved historical ancestry must remain a stated limitation even if the new cohort passes contemporary overlap checks.

The intended evaluation reports technical scoring coverage on all eligible ECGs, class-specific known reference counts, sensitivity, specificity, predictive values, F1, discrimination and calibration where both reference polarities support them. Patient or subject/signal components should define resampling units when repeated recordings are present. Any abstention policy requires its own frozen calibration cohort and a ledger separating emitted decisions, withheld known positives and technical failures. No independent-study sample size, completed inference count or performance result is asserted here. The current manuscript reports only the completed retrospective analyses and leaves independent external evaluation open.
